\editionPage[\pdfbookmark[0]{Corrigés des bilans}{bm-bilan-A-cor-auto-parent}]{bilan-A-cor-auto}{Bilan A · Corrigé}{Automatismes expliqués}{\hyperref[bilan-A-auto]{Énoncé p.\,\pageref*{bilan-A-auto}}}
\begin{blocEdition}{bilan-A-cor-auto}
\begin{questions}
\item $\res{3\,040\,006}$.
\item $623-100+2=\res{525}$.
\item $4\times7=28$ puis trois zéros : $\res{28\,000}$.
\item $40\times5=\res{200}$.
\item $\res{1, 2, 4, 5, 10, 20}$.
\item $\res{3{,}25}$.
\item $\res{3+\frac25}$.
\item $\frac68+\frac18=\res{\frac78}$.
\item $27:3\times2=\res{18}$.
\item $5{,}090<5{,}100$ : $\res{5{,}09<5{,}1}$.
\item $\res{456}$.
\item $\res{4{,}7}$.\end{questions}
\end{blocEdition}
\editionPage{bilan-A-cor-1}{Bilan A · Corrigé 1}{Le concert de la région}{\hyperref[bilan-A-ex-1]{Énoncé p.\,\pageref*{bilan-A-ex-1}}}
\begin{blocEdition}{bilan-A-cor-1}
\editionEtape{Total}$18\,750+16\,480=\res{35\,230}$ spectateurs.\editionEtape{Recette}$18\,750\times24=\res{450\,000}$ euros ; estimation : $19\,000\times24\approx456\,000$.\editionEtape{Places vides}Capacité : $40\,000$ places ; vides : $40\,000-35\,230=\res{4\,770}$.
\end{blocEdition}
\editionPage{bilan-A-cor-2}{Bilan A · Corrigé 2}{Les sachets de bonbons}{\hyperref[bilan-A-ex-2]{Énoncé p.\,\pageref*{bilan-A-ex-2}}}
\begin{blocEdition}{bilan-A-cor-2}
\editionEtape{Sachets}$1\,000=8\times125$ : $\res{125}$ sachets, aucun bonbon restant.\editionEtape{Recette}$125\times2{,}50=\res{312{,}50}$ euros.\editionEtape{Partage}$45:4=\res{11{,}25}$ euros chacun ; contrôle : $11{,}25\times4=45$.
\end{blocEdition}
\editionPage{bilan-A-cor-3}{Bilan A · Corrigé 3}{La pizza partagée}{\hyperref[bilan-A-ex-3]{Énoncé p.\,\pageref*{bilan-A-ex-3}}}
\begin{blocEdition}{bilan-A-cor-3}
\editionEtape{Ensemble}$\frac38+\frac28=\res{\frac58}$.\editionEtape{Reste}$1-\frac58=\res{\frac38}$.\editionEtape{Masse}$480:8\times3=\res{180}$ g.
\end{blocEdition}
\editionPage{bilan-A-cor-4}{Bilan A · Corrigé 4}{Le programme mystère}{\hyperref[bilan-A-ex-4]{Énoncé p.\,\pageref*{bilan-A-ex-4}}}
\begin{blocEdition}{bilan-A-cor-4}
\editionEtape{Avec 4}$4\times3=12$ ; $12+6=18$ ; $18\times2=\res{36}$.\editionEtape{Remonter}$48:2=24$ ; $24-6=18$ ; $18:3=\res6$.\editionEtape{Pour N}Le résultat s’écrit $\res{((N\times3)+6)\times2}$.
\end{blocEdition}
\editionPage{bilan-B-cor-auto}{Bilan B · Corrigé}{Automatismes expliqués}{\hyperref[bilan-B-auto]{Énoncé p.\,\pageref*{bilan-B-auto}}}
\begin{blocEdition}{bilan-B-cor-auto}
\begin{questions}
\item $\res{2\,500}$ m.
\item $\res{0{,}75}$ kg.
\item $\res{300\ \mathrm{cm}^2}$.
\item $\res{54\ \mathrm{cm}^2}$.
\item $\res{150}$ s.
\item $\res{45^\circ}$.
\item $\res6$.
\item $\res{12}$.
\item Trois fois plus : $\res{18}$ euros.
\item $\res{2\ \text{chances sur }6}$.
\item $\res7$ cm.
\item $\res{18}$ cm.\end{questions}
\end{blocEdition}
\editionPage{bilan-B-cor-1}{Bilan B · Corrigé 1}{La rénovation de la chambre}{\hyperref[bilan-B-ex-1]{Énoncé p.\,\pageref*{bilan-B-ex-1}}}
\begin{blocEdition}{bilan-B-cor-1}
\editionEtape{Aire}$4{,}5\times3{,}2=\res{14{,}4\ \mathrm m^2}$.\editionEtape{Prix}$14{,}4\times25=\res{360}$ euros.\editionEtape{Plinthe}Périmètre : $2\times(4{,}5+3{,}2)=15{,}4$ m ; sans la porte : $15{,}4-0{,}8=\res{14{,}6}$ m.
\end{blocEdition}
\editionPage{bilan-B-cor-2}{Bilan B · Corrigé 2}{Le marathon des enfants}{\hyperref[bilan-B-ex-2]{Énoncé p.\,\pageref*{bilan-B-ex-2}}}
\begin{blocEdition}{bilan-B-cor-2}
\editionEtape{Le premier}$9$ h 45 min 30 s $+38$ min 45 s : 30 s mènent à 9 h 46, puis 38 min 15 s : arrivée à $\res{10\ \text h\ 24\ \text{min}\ 15\ \text s}$.\editionEtape{Le dernier}De 9 h 45 min 30 s à 10 h 41 : $\res{55\ \text{min}\ 30\ \text s}$.\editionEtape{Écart}$55$ min 30 s $-38$ min 45 s $=\res{16\ \text{min}\ 45\ \text s}$.
\end{blocEdition}
\editionPage{bilan-B-cor-3}{Bilan B · Corrigé 3}{Le sondage de la classe}{\hyperref[bilan-B-ex-3]{Énoncé p.\,\pageref*{bilan-B-ex-3}}}
\begin{blocEdition}{bilan-B-cor-3}
\editionEtape{Effectifs}À pied : $28:2=\res{14}$ ; bus : $28:4=\res7$ ; vélo : $28-14-7=\res7$.\editionEtape{Secteur}$\frac7{28}=\res{\frac14}$ du disque.\editionEtape{Tirage}Les 28 fiches sont équiprobables : $\res{7\ \text{chances sur }28}$, soit une chance sur quatre.
\end{blocEdition}
\editionPage{bilan-B-cor-4}{Bilan B · Corrigé 4}{La figure à construire}{\hyperref[bilan-B-ex-4]{Énoncé p.\,\pageref*{bilan-B-ex-4}}}
\begin{blocEdition}{bilan-B-cor-4}
\editionEtape{Le carré}Périmètre $4\times6=\res{24}$ cm ; aire $6\times6=\res{36\ \mathrm{cm}^2}$.\editionEtape{Le cercle}Le rayon 3 cm est la moitié de AB : le cercle coupe [AB] en son $\res{\text{milieu}}$.\editionEtape{L’axe}La diagonale $\res{(AC)}$ est un axe de symétrie du carré qui passe par A.
\end{blocEdition}
